% Generated draft: requires mathematical and source review.
\documentclass[11pt]{article}
\usepackage[margin=25mm]{geometry}
\usepackage{fontspec}
\setmainfont{Latin Modern Roman}
\usepackage{amsmath,amssymb,mathtools,unicode-math}
\setmathfont{Latin Modern Math}
\usepackage{enumitem,needspace,xurl,hyperref}
\setlength{\parindent}{0pt}
\setlength{\parskip}{5pt}
\setlength{\emergencystretch}{3em}
\newcommand{\N}{\mathbb{N}}
\newcommand{\Z}{\mathbb{Z}}
\newcommand{\Q}{\mathbb{Q}}
\newcommand{\R}{\mathbb{R}}
\newcommand{\C}{\mathbb{C}}
\newcommand{\F}{\mathbb{F}}
\newcommand{\A}{\mathbb{A}}
\newcommand{\PP}{\mathbb{P}}
\newcommand{\E}{\mathbb{E}}
\newcommand{\eps}{\varepsilon}
\newcommand{\dd}{\,\mathrm{d}}
\begin{document}
\section*{UnsolvedMath ID: 30004235}
\textbf{Source problem code:} catalogue entry
\subsection{Definitions and mathematical statement}
Let \(n\) be a positive integer, write \([n]=\{1,\ldots,n\}\), and let
\[
V_n=\R^{\binom n2}=\{(x_{ab})_{1\leq a<b\leq n}:x_{ab}\in\R\}.
\]
The source indexes its hyperplanes by unordered triples but does not specify how the indices in its equation are ordered. The following is therefore a conditional formulation using increasing representatives; this convention requires confirmation against the source:
\[
H_{ijk}=\{x\in V_n:x_{ij}+x_{jk}-x_{ik}=0\},\qquad
\mathcal H_n=\{H_{ijk}:1\leq i<j<k\leq n\}.
\]
Under this convention, \(\mathcal H_n\) is empty when \(n<3\). The source does not explicitly specify the range of \(n\); positive integers are used here to state the family.

Define the intersection lattice by
\[
L(\mathcal H_n)=\left\{\bigcap_{H\in S}H:S\subseteq\mathcal H_n\right\},
\qquad \bigcap_{H\in\varnothing}H=V_n,
\]
ordered by reverse inclusion. For \(X\in L(\mathcal H_n)\), define the integers \(\mu(V_n,X)\) by \(\mu(V_n,V_n)=1\) and
\[
\mu(V_n,X)=-\sum_{\substack{Y\in L(\mathcal H_n)\\ V_n\leq Y<X}}\mu(V_n,Y)
\quad\text{when }X\ne V_n.
\]
Use the ambient-space convention for the characteristic polynomial:
\[
\chi_{\mathcal H_n}(t)=\sum_{X\in L(\mathcal H_n)}\mu(V_n,X)t^{\dim X}.
\]
A region is a connected component of
\[
V_n\setminus\bigcup_{H\in\mathcal H_n}H.
\]

\textbf{Requested determination.} For each \(n\), describe \(L(\mathcal H_n)\), determine \(\chi_{\mathcal H_n}(t)\), and describe and enumerate the regions. The supplied cleaned statement also requests other combinatorial invariants without specifying which ones. Its original wording asks more broadly what is known about the arrangement. Neither wording supplies a particular conjectural formula or theorem to prove. All these requests remain conditional on resolving the indexing convention above.
\subsection{Short English statement}
Study the real hyperplane arrangement defined by triangle equations among coordinates indexed by pairs of vertices. Describe its intersection lattice, characteristic polynomial, and regions as the number of vertices varies. The draft uses increasing vertex indices explicitly because the supplied unordered-triple notation leaves the equation's ordering unclear. The supplied open label and literature assessment are dated curation, not independently verified conclusions.
\subsection{Sources}
\begin{itemize}
\item Supplied local UnsolvedMath registry record, permanent ID \(30004235\), source problem code \texttt{catalogue entry}, title \emph{Triangle Equation Hyperplane Arrangements}; collection: Oberwolfach Reports Open Problems. Registry locations identified by the project are \texttt{data/unsolvedmath/problems.json} and the supplementary \texttt{catalog/problems.json}; the supplied record is the evidence used here.
\item Supplied report citation: \emph{Geometric, Algebraic, and Topological Combinatorics (2020)}, Oberwolfach Reports, pp.~2395--2472, DOI \texttt{10.4171/owr/2019/39}. \url{https://doi.org/10.4171/owr/2019/39}. The record gives proposed year 2019 and identifies no proposer. Its background describes the reports as records of mathematical workshops and their open-problem discussions.
\item Supplied literature reference, labeled ``Linial arrangement literature'': \url{https://doi.org/10.1006/jcta.2000.3106}. No author, title, or specific supporting result is supplied for this reference. The supplied assessment distinguishes the triangle-equation arrangement from classical Linial arrangements.
\item Literature assessment dated exactly \texttt{2026-08-22}: the supplied curation labels the problem open, states that classical Linial arrangements are well studied, and reports no identified complete treatment of the distinct triangle-equation arrangement, with its requested properties remaining open. Its listed sources are the Oberwolfach report and the Linial arrangement reference above. This is supplied dated curation, not independently verified; the record itself requests independent verification of time-sensitive status and cited claims.
\item The supplied statement status is \texttt{corrected\_verified}. Its verification note says that the official report was checked directly and that the cleaned statement preserves the mathematical target while restoring definitions and removing source-dependent wording. This is a supplied verification claim; no independent source visit or mathematical verification is asserted here.
\end{itemize}
\end{document}
